%%%%%%%%%%%%%%%%%%%%%%%%%%%%%%%%%%%%%%%%%%%%%%%%%%%%%%%%%%%%%%%%%%%%%%%%%%%%%%%%%%%%%%%%%%%%%%%%%%%%%%%%%%%%%%%%%%%%%%%%%%%%%%%%%%%%%%%%%%%%%%%%%%%%%%%%%%%
% This is just an example/guide for you to refer to when submitting manuscripts to Frontiers, it is not mandatory to use Frontiers .cls files nor frontiers.tex  %
% This will only generate the Manuscript, the final article will be typeset by Frontiers after acceptance.   
%                                              %
%                                                                                                                                                         %
% When submitting your files, remember to upload this *tex file, the pdf generated with it, the *bib file (if bibliography is not within the *tex) and all the figures.
%%%%%%%%%%%%%%%%%%%%%%%%%%%%%%%%%%%%%%%%%%%%%%%%%%%%%%%%%%%%%%%%%%%%%%%%%%%%%%%%%%%%%%%%%%%%%%%%%%%%%%%%%%%%%%%%%%%%%%%%%%%%%%%%%%%%%%%%%%%%%%%%%%%%%%%%%%%

%%% Version 3.4 Generated 2018/06/15 %%%
%%% You will need to have the following packages installed: datetime, fmtcount, etoolbox, fcprefix, which are normally inlcuded in WinEdt. %%%
%%% In http://www.ctan.org/ you can find the packages and how to install them, if necessary. %%%
%%%  NB logo1.jpg is required in the path in order to correctly compile front page header %%%

\documentclass[utf8]{FrontiersinHarvard} % for articles in journals using the Harvard Referencing Style (Author-Date), for Frontiers Reference Styles by Journal: https://zendesk.frontiersin.org/hc/en-us/articles/360017860337-Frontiers-Reference-Styles-by-Journal
%\documentclass[utf8]{FrontiersinVancouver} % for articles in journals using the Vancouver Reference Style (Numbered), for Frontiers Reference Styles by Journal: https://zendesk.frontiersin.org/hc/en-us/articles/360017860337-Frontiers-Reference-Styles-by-Journal
%\documentclass[utf8]{frontiersinFPHY_FAMS} % Vancouver Reference Style (Numbered) for articles in the journals "Frontiers in Physics" and "Frontiers in Applied Mathematics and Statistics" 

%\setcitestyle{square} % for articles in the journals "Frontiers in Physics" and "Frontiers in Applied Mathematics and Statistics" 
\usepackage{url,hyperref,lineno,microtype,subcaption}
\usepackage[onehalfspacing]{setspace}

\linenumbers


% Leave a blank line between paragraphs instead of using \\


\def\keyFont{\fontsize{8}{11}\helveticabold }
\def\firstAuthorLast{Sample {et~al.}} %use et al only if is more than 1 author
\def\Authors{Etienne Boileau\,$^{1,2,3}$, Christoph Dieterich\,$^{1,2,3}$}
% Affiliations should be keyed to the author's name with superscript numbers and be listed as follows: Laboratory, Institute, Department, Organization, City, State abbreviation (USA, Canada, Australia), and Country (without detailed address information such as city zip codes or street names).
% If one of the authors has a change of address, list the new address below the correspondence details using a superscript symbol and use the same symbol to indicate the author in the author list.
\def\Address{$^{1}$Section of Bioinformatics and Systems Cardiology, Klaus Tschira Institute for Integrative Computational Cardiology, Im Neuenheimer Feld 669, 69120 Heidelberg, Germany
\\
$^{2}$Department of Internal Medicine III (Cardiology, Angiology, and Pneumology), University Hospital Heidelberg, Im Neuenheimer Feld 669, 69120 Heidelberg, Germany
\\
$^{3}$DZHK (German Centre for Cardiovascular Research) Partner Site Heidelberg/Mannheim}
% The Corresponding Author should be marked with an asterisk
% Provide the exact contact address (this time including street name and city zip code) and email of the corresponding author
\def\corrAuthor{Christoph Dieterich}

\def\corrEmail{christoph.dieterich@uni-heidelberg.de}


%%%%%%%%%%%%%%%%%%%%%%%%%%%%%%%%%%%%%%%%%%%%%%%%%%%%%%%
\graphicspath{{figures/}}

\usepackage{xspace}

\newcommand{\insi}{\textit{in situ}\xspace}
\newcommand{\exsi}{\textit{ex situ}\xspace}
\newcommand{\apri}{\textit{a priori}\xspace}
\newcommand{\vivo}{\textit{in vivo}\xspace}
\newcommand{\vitro}{\textit{in vitro}\xspace}
\newcommand{\denovo}{\textit{de novo}\xspace}
\newcommand{\ie}{\textit{i.e.}\xspace}
\newcommand{\eg}{\textit{e.g.}\xspace}
\newcommand{\etc}{\textit{etc.}\xspace}
\newcommand{\vs}{\textit{vs.}\xspace}
\newcommand{\etl}{\textit{et al.}\xspace}
\newcommand{\perse}{\textit{per se}\xspace}

\newcommand{\rnaseq}{\textsc{RNA}-seq\xspace}
\newcommand{\scn}{\textsc{scNapBar}\xspace}
\newcommand{\scnast}{\textsc{scNaST}\xspace}

% revision
\usepackage{xcolor}
\newcommand*{\red}{\textcolor{red}}

%%%%%%%%%%%%%%%%%%%%%%%%%%%%%%%%%%%%%%%%%%%%%%%%%%%%%%%


\begin{document}
\onecolumn
\firstpage{1}

\title[Running Title]{Article Title} 

\author[\firstAuthorLast ]{\Authors} %This field will be automatically populated
\address{} %This field will be automatically populated
\correspondance{} %This field will be automatically populated

\extraAuth{}% If there are more than 1 corresponding author, comment this line and uncomment the next one.
%\extraAuth{corresponding Author2 \\ Laboratory X2, Institute X2, Department X2, Organization X2, Street X2, City X2 , State XX2 (only USA, Canada and Australia), Zip Code2, X2 Country X2, email2@uni2.edu}


\maketitle


\begin{abstract}

%%% Leave the Abstract empty if your article does not require one, please see the Summary Table for full details.
\section{}

We introduce Single-cell Nanopore Spatial Transcriptomics (\scnast), a set of tools to facilitate the analysis of spatial gene expression from second- and third-generation sequencing, allowing to generate a full-length single-cell transcriptional landscape of the tissue microenvironment.
%short and long read sequencing of cDNA synthesized \insi on tissue sections
Taking advantage of the Visium Spatial platform, we adapted a strategy recently developed to assign barcodes to long-read single-cell sequencing data for spatial capture technology.



\tiny
 \keyFont{ \section{Keywords:} Spatial transcriptomics, Single-cell RNA sequencing, Oxford Nanopore, keyword, keyword} 
\end{abstract}

\section*{Introduction}

%For Original Research Articles the introduction should be succinct, with no subheadings 


% For Original Research articles, please note that the Material and Methods section can be placed in any of the following ways: before Results, before Discussion or after Discussion.

\section*{Material and Methods}


\subsection*{Mouse heart samples}


\subsection*{10X Genomics Visium experiments}


\subsection*{Oxford Nanopore sequencing libraries}


\subsection*{Cell barcode assignment}



\section*{Results}
Fresh-frozen tissue samples were stained, imaged and fixed on Visium Spatial Gene Expression Slides (10X Genomics) for permeabilization and \insi RNA capture.
Full-length cDNA libraries were split for the preparation of 3' short-read and long-read Nanopore sequencing libraries.
Short-read data were used for the assignment of spatial barcodes to Nanopore reads using the \scn workflow, and subsequently used to define anatomical regions within the tissue organization (Fig.~\ref{fig:1}a).
Long-read data were used for transcriptome assembly and transcript abundance quantification, and layered onto the stained images to reveal the spatial organization of isoform expression.


% >>> CARDIAC RELEVANCE
% background cardiac spatial transcriptomics, relevance to the field, etc.
% We demonstrate WHAT in 4 heart slice after MI?



\section*{Discussion}



\section*{Data Availability}
The datasets [GENERATED/ANALYZED] for this study can be found in the [NAME OF REPOSITORY] [LINK].
% Please see the availability of data guidelines for more information, at https://www.frontiersin.org/about/author-guidelines#AvailabilityofData

\section*{Code Availability}


\section*{Author Contributions}
CD supervised the research.
EB analyzed the data and wrote the manuscript. 
XX performed the experiments.
All authors contributed to review and editing.


\section*{Funding}
\red{Is this still correct?}
EB and CD acknowledge support by the Klaus Tschira Stiftung gGmbH [00.219.2013]. 
CD acknowledge the DZHK (German Centre for Cardiovascular Research) Partner Site Heidelberg/Mannhein.

%\section*{Acknowledgments}

\section*{Conflict of Interest}
The authors declare that they have no conflict of interest.


\section*{Supplemental Data}
% \href{http://home.frontiersin.org/about/author-guidelines#SupplementaryMaterial}{Supplementary Material} should be uploaded separately on submission, if there are Supplementary Figures, please include the caption in the same file as the figure. LaTeX Supplementary Material templates can be found in the Frontiers LaTeX folder.


\bibliographystyle{Frontiers-Harvard} %  Many Frontiers journals use the Harvard referencing system (Author-date), to find the style and resources for the journal you are submitting to: https://zendesk.frontiersin.org/hc/en-us/articles/360017860337-Frontiers-Reference-Styles-by-Journal. For Humanities and Social Sciences articles please include page numbers in the in-text citations 
%\bibliographystyle{Frontiers-Vancouver} % Many Frontiers journals use the numbered referencing system, to find the style and resources for the journal you are submitting to: https://zendesk.frontiersin.org/hc/en-us/articles/360017860337-Frontiers-Reference-Styles-by-Journal
\bibliography{test}

%%% Make sure to upload the bib file along with the tex file and PDF
%%% Please see the test.bib file for some examples of references

\section*{Figure captions}

%%% Please be aware that for original research articles we only permit a combined number of 15 figures and tables, one figure with multiple subfigures will count as only one figure.
%%% Use this if adding the figures directly in the mansucript, if so, please remember to also upload the files when submitting your article
%%% There is no need for adding the file termination, as long as you indicate where the file is saved. In the examples below the files (logo1.eps and logos.eps) are in the Frontiers LaTeX folder
%%% If using *.tif files convert them to .jpg or .png
%%%  NB logo1.eps is required in the path in order to correctly compile front page header %%%

\begin{figure}[h!]
\begin{center}
\includegraphics[width=10cm]{logo1}% This is a *.eps file
\end{center}
\caption{ Enter the caption for your figure here.  Repeat as  necessary for each of your figures}\label{fig:1}
\end{figure}


\begin{figure}[h!]
\begin{center}
\includegraphics[width=15cm]{logos}
\end{center}
\caption{This is a figure with sub figures, \textbf{(A)} is one logo, \textbf{(B)} is a different logo.}\label{fig:2}
\end{figure}

%%% If you are submitting a figure with subfigures please combine these into one image file with part labels integrated.
%%% If you don't add the figures in the LaTeX files, please upload them when submitting the article.
%%% Frontiers will add the figures at the end of the provisional pdf automatically
%%% The use of LaTeX coding to draw Diagrams/Figures/Structures should be avoided. They should be external callouts including graphics.

\end{document}
